\subsection{Translations and uniform Lipschitz regularity}

Extend $u_D,u_A$ by one outside their intervals. The extensions are continuous at the endpoints because the endpoint fibers have length
zero. For $s\in\mathbb R$ define \[
 \widehat r_R(s)=\Bigl[x+b\int_0^1r_C(c)\,\widehat u_D(s-c)\,\widehat u_A(s+c)\,dc\Bigr]^{-1}.
\] Since $\int r_C=1$, the denominator lies between $b_-\kappa^2$ and $1+b_+$, so $\widehat r_R\in[m_0,\overline K]$, and the extension
agrees with $r_R$ on $[0,1]$. For a translation $t$, let $r_R^t(r)=\widehat r_R(r+t)$, $r_C^t=r_C$, $u_D^t(d)=\widehat u_D(d+t)$,
$u_A^t(s)=\widehat u_A(s+t)$, and $v^t=\log(r_R^t,r_C^t,u_D^t,u_A^t)$. With $e^t=\mathcal G_{x,b}(v^t)$ we have \[
 e_R^t=0,\qquad |e_C^t|,|e_D^t|,|e_A^t|\le2\overline B|t|,\qquad
 \|e^t\|_{\mathcal H}\le2\sqrt5\,\overline B|t| .
\] Indeed, the definition of $\widehat r_R$ gives the row equation exactly. In the column equation, change variables from $r$ to $r+t$.
The domains differ by two strips of total length at most $2|t|$, on which the extended density is at most $\overline B$. For the
diagonals and antidiagonals, make the same change of variable. The original domain of the translated diagonal differs by the same
strips. Outside the diagonal interval the original fiber is empty and the extended vacancy equals one. The bounds hold along every
segment. Since $v_C^t=v_C$, Lemma~\ref{lem:coercividad} gives $\|v^t-v\|_{\mathcal H}\le2\sqrt5\,\overline B|t|/\lambda_c$.

\emph{From the $L^2$ norm to the uniform norm.} Along a fiber of $f$, each coordinate of the other families varies with a slope of
absolute value one or two. Telescoping three factors bounded by $\overline K$ and applying Cauchy--Schwarz and
$|e^\alpha-e^\beta|\le\overline K|\alpha-\beta|$ gives $|J_f(v^t)-J_f(v)|\le\sqrt3\,\overline K^3\|v^t-v\|_{\mathcal H}$. The
denominators in the row and column equations are $x+bJ_f\ge b_-m_0\kappa^2$, and those in the diagonal equations are at least one. With
\[
 d_0=\min\{1,b_-m_0\kappa^2\},\qquad
 L_*=\frac{2\overline B}{d_0}+\frac{2\sqrt{15}\,b_+\overline K^3\overline B}{d_0^2\lambda_c},
\] and writing $y_f$ for the corresponding normalized factor, we have $y_f^t(d_f+bJ_f(v^t))=1+e_f^t$ and $y_f(d_f+bJ_f(v))=1$, with
$d_{R,C}=x$ and $d_{D,A}=1$. Comparing reciprocals gives $|y_f^t-y_f|\le L_*|t|$. The row translation controls $R$, $D$, and $A$. The
symmetric construction, which translates $C$ and keeps $R$ fixed, controls $C$.

\begin{teorema}[Uniform regularity up to the full board, {\cite[Theorem~7.7]{Perfil}}]\label{thm:lipschitz} The optimal
profiles~\cite[Theorem~3.1]{Perfil}, rescaled and normalized by $\int r_R=\int r_C=1$, extend continuously to $x=0$ and satisfy, with
constants independent of $x\in[0,1/2]$, \[
 \|r_R\|_\infty,\|r_C\|_\infty\le\overline K,\quad \|u_D\|_\infty,\|u_A\|_\infty\le1,\quad
 \Lip(r_R,r_C,u_D,u_A)\le L_*,
\] \[
 \|\nu\|_\infty\le\overline B,\quad \Lip_\infty\nu\le6b_+\overline K^2L_*,\quad
 \|\Theta_x\|_\infty\le\overline K^2,\quad \Lip_\infty\Theta_x\le2\overline KL_*,
\] where $\Theta_x(r,c)=r_R(r)r_C(c)$. \end{teorema}

\begin{proof} Continuity at the endpoint is Proposition~5.3 of~\cite{Perfil}. The preceding proof applies at the endpoint as well.
Alternatively, the Lipschitz bounds pass to the limit by uniform convergence. For $\nu$, telescoping the four factors contributes at
most $b_+\overline K^2L_*(|\Delta r|+|\Delta c|+|\Delta r-\Delta c|+|\Delta r+\Delta c|)
\le6b_+\overline K^2L_*\|\Delta(r,c)\|_\infty$, and the product of the two factors of $\Theta_x$ contributes $2\overline KL_*$. \end{proof}

Continuity in the parameter $x$ is uniform in space (\cite[Propositions~5.3 and~7.5]{Perfil}), and \cite[Theorem~5.4]{Perfil} even
gives an analytic branch near zero in the space of continuous functions. We claim no spatial $C^1$ regularity. The bounds for the
continuous problem do not by themselves prove any discrete estimate.
